\section{Implementación}
\label{sec:implementacion}

Los dos modelos se implementaron en C++20 como una biblioteca, \texttt{refor\_core}, sobre la que se construyen tres programas: \texttt{refor}, una herramienta de línea de comandos para análisis y reportes; \texttt{refor\_ple}, que resuelve por separado el modelo entero de la \cref{subsec:espacial-entera}; y \texttt{refor\_gui}, una aplicación de escritorio para los técnicos de planeación. Los programas enteros se resuelven con HiGHS. El código, los datos procesados y las instrucciones para reproducir cada resultado de este trabajo están en el repositorio del proyecto.

\subsection{Organización}
\label{subsec:impl-organizacion}

La \cref{tab:modulos} relaciona cada módulo de la biblioteca con la parte del modelo que implementa. El flujo sigue la descomposición de las \cref{sec:modelo-compra,sec:modelo-espacial}: se leen los datos de INECOL, se calculan las curvas $\hat{x}_s^{*}(n)$ y $\hat{D}_s(n)$ de cada especie y polígono, se resuelve la mochila de cada polígono, se reparte su compra en un mosaico de 52 posiciones y se agrupan los polígonos con la misma composición, de modo que cada mosaico distinto se optimiza una sola vez.

\begin{table}[t]
\centering
\caption{Módulos de \texttt{refor\_core} y la parte del modelo que implementan.}
\label{tab:modulos}
\small
\begin{tabularx}{\columnwidth}{@{}>{\RaggedRight\arraybackslash}p{2.6cm}>{\RaggedRight\arraybackslash}X@{}}
\toprule
Módulo & Contenido \\
\midrule
\texttt{simulacion} & Escenarios Monte Carlo y curvas $\hat{x}^{*}(n)$, $\hat{D}(n)$ \\
\texttt{exacto} & Las mismas curvas con la cola binomial exacta \\
\texttt{mochila} & Programación dinámica y programa entero de la mochila de elección múltiple \\
\texttt{modelo} & Curvas en paralelo y solución de cada polígono \\
\texttt{plan} & Composición por mayor residuo, agrupación de polígonos y corridas del memético \\
\texttt{memetico/} & Mosaico toroidal, matriz $W$, búsqueda tabú y algoritmo memético \\
\texttt{asignacion\_exacta/} & Modelo entero del mosaico y comparación con el memético \\
\bottomrule
\end{tabularx}
\end{table}

\subsection{Reproducibilidad}
\label{subsec:impl-reproducibilidad}

Todo resultado depende de una sola semilla (42 por defecto) y debe ser idéntico sin importar el número de hilos ni la plataforma. Para lograrlo, cada flujo aleatorio se inicializa a partir de su identidad lógica y no del orden en que se ejecuta: los escenarios de la especie $s$ en el polígono $i$ usan una semilla derivada de $(42,i,s)$, y cada corrida del memético, una derivada del mosaico y del número de corrida. El generador es \texttt{mt19937\_64}, cuya secuencia fija el estándar de C++. En cambio, los algoritmos de las distribuciones de la biblioteca estándar difieren entre implementaciones, por lo que los enteros uniformes y las permutaciones se generan con funciones propias, y la supervivencia de cada planta se decide comparando la salida del generador con $q_s\cdot2^{64}$. La biblioteca se compila sin contracción de operaciones de punto flotante, que podría cambiar el redondeo entre compiladores. Como comprobación, la compra de los 30 polígonos calculada con un hilo y con diez da archivos idénticos byte por byte. El flujo de publicación compila los programas en Linux, macOS y Windows y compara sus resultados.

\subsection{Pruebas y desempeño}
\label{subsec:impl-pruebas}

La biblioteca tiene 30 casos de prueba automatizados. Los más importantes comparan cada método con una referencia independiente: la cola binomial con valores conocidos; las curvas de Monte Carlo con las exactas; la programación dinámica y el programa entero de la mochila con la enumeración exhaustiva de metas, con y sin presupuesto; y el modelo entero del mosaico con la enumeración exhaustiva en un mosaico de $3\times3$. Otras pruebas verifican invariantes: la monotonía de $\hat{x}^{*}(n)$ en $n$ y en $\alpha$, que la búsqueda tabú nunca empeora su solución inicial, que el cruce conserva la composición, que la distancia entre mosaicos reconoce un mismo esquema desplazado y que la secuencia aleatoria del memético es la misma en todas las plataformas.

En una computadora con un procesador Apple M4 de diez núcleos, las curvas de los 30 polígonos con $K=10{,}000$ escenarios se calculan en 1.8~s, y en 0.15~s con la cola binomial exacta. Las 30 corridas del memético para cada uno de los seis mosaicos distintos toman en conjunto unos 30~s. El modelo entero es el único paso costoso: se le dan 300~s por mosaico.
